% ============================================================
% BÖLÜM 1 - GİRİŞ
% ============================================================
\chapter{GİRİŞ}

\section{Problem Tanımı}

Akışkanlar mekaniği, enerji ve ısı transferi süreçlerinin temelini oluşturan disiplinlerden biridir; binaların ısıtma ve havalandırması, elektronik bileşenlerin soğutulması, atmosfer ve okyanus dinamiği, jeotermal sistemler ve nükleer reaktör çekirdek soğutması gibi son derece farklı uygulama alanlarında karşımıza çıkar. Bu uygulamaların büyük bir kısmında akış yalnızca atalet ile değil, aynı zamanda sıcaklık farkından doğan kaldırma kuvveti ile sürüklenir. Atalet kuvvetlerinin (zorlanmış konveksiyon) ve kaldırma kuvvetlerinin (doğal konveksiyon) birlikte etkili olduğu bu rejim, akışkanlar mekaniğinde \emph{karışık konveksiyon} (\textit{mixed convection}) olarak adlandırılır \cite{bejan_convection_2013, incropera_heat_transfer_2017}.

Karışık konveksiyon, Reynolds sayısının yanı sıra Rayleigh sayısı ve Prandtl sayısının da birlikte rol oynadığı çok ölçekli bir problemdir. Rejimin baskın karakteri Richardson sayısı $\mathrm{Ri}=\mathrm{Ra}/(\mathrm{Re}^{2}\mathrm{Pr})$ ile karakterize edilir: $\mathrm{Ri}\ll 1$ atalet baskın (zorlanmış konveksiyon), $\mathrm{Ri}\gg 1$ kaldırma baskın (doğal konveksiyon), $\mathrm{Ri}\sim\mathcal{O}(1)$ ise iki etkinin karşılaştırılabilir düzeyde olduğu \emph{gerçek} karışık konveksiyon rejimidir. Bu son rejim, hem büyük ölçekli düzenli yapıların (örneğin termal pleymler) hem de küçük ölçekli türbülanslı dalgalanmaların aynı anda var olduğu, dolayısıyla hem büyük girdap dinamiğinin hem de alt-ızgara modellemesinin önemli olduğu zorlu bir test problemidir.

Üç boyutlu sıkıştırılamaz akış denklemleri (Navier--Stokes), enerji denklemi ve Boussinesq yaklaşımı altında kaldırma kuvveti terimi birlikte ele alındığında ortaya çıkan denklemler, doğrusal olmayan adveksiyon terimi nedeniyle analitik çözüme kapalıdır. Klasik sayısal yöntemler doğrudan sayısal simülasyon (DNS), büyük girdap simülasyonu (LES) ve Reynolds ortalaması tabanlı yöntemler (RANS) çerçevesinde toplanır. DNS, $\mathrm{Re}=10\,000$ mertebesinde dahi $\mathcal{O}(10^{9})$ ızgara noktası ve son derece küçük zaman adımları gerektirdiği için pratik mühendislik problemlerinde uygulanabilir değildir \cite{pope_turbulence_2000}. LES ise enerji içeren büyük ölçekleri açıkça çözer, küçük ölçekli yapıları ise alt-ızgara modeli ile yaklaşık olarak temsil eder; bu sayede DNS'ye kıyasla çok daha az kaynakla çözülebilir, ancak hâlâ yüksek hesaplama maliyeti içerir.

Son yıllarda, türbülanslı akışların modellenmesinde makine öğrenmesi tabanlı yaklaşımların hızla geliştiğine tanık olunmaktadır. Bu yaklaşımlar üç ana eksende incelenebilir: (i) saf veri tabanlı operatör öğrenme yöntemleri (Fourier neural operator, DeepONet) \cite{li_fno_2021, lu_deeponet_2021}; (ii) fiziksel kısıtları kayıp fonksiyonuna yerleştiren fizik bilgili sinir ağları (PINN) \cite{raissi_pinn_2019}; (iii) klasik bir sayısal çözücüye küçük ölçekli düzeltmeler öğreten hibrit yaklaşımlar \cite{kochkov_ml_cfd_2021}. Bu üç eksen birbirini dışlamamakla birlikte, her birinin üç boyutlu türbülans bağlamında belirli zayıf yönleri vardır. Saf veri tabanlı yöntemler büyük etiketli veri kümelerine ihtiyaç duyar; PINN yaklaşımları yüksek mertebe türevlerin otomatik türev üzerinden hesaplanmasından kaynaklanan kararsızlık ve kayıp dengesi sorunları yaşar \cite{krishnapriyan_failure_2021, wang_gradient_2021}; klasik çözücüye dayalı hibrit yöntemler ise bağımlı oldukları çözücünün karmaşıklığını miras alır.

Bu çalışmanın ön döneminde (\textit{bitirme projesi 1}) bu üç yaklaşımdan PINN paradigması, üç boyutlu Taylor--Green girdabı (TGV3D) probleminde ayrıntılı olarak incelenmiştir \cite{tezgocen_bitirme1_2026}. Söz konusu çalışmada yapısal başlangıç koşulu ansatzı, çok dallı aktivasyon mimarisi, müfredat öğrenmesi ve fiziksel kısıtların kademeli aktivasyonu gibi tekniklerle PINN eğitimi başarılı şekilde stabilize edilmiş; ancak yüksek Reynolds sayılarında küçük ölçekli yapıların temsilinde ve uzun zaman ufuklarında enerji dengesinin korunmasında temel sınırlamalar gözlemlenmiştir. Bu sınırlamalar, fiziksel kısıtların \emph{ceza} biçiminde değil, ağın \emph{yapısı} içinde temsil edildiği alternatif bir paradigmanın araştırılmasını motive etmiş ve INNATE (\textit{Intrinsic Navier--Stokes Neural Architecture for Temporal Evolution}) yaklaşımı önerilmiştir.

INNATE'in temel ilkesi, fiziksel operatörlerin --- adveksiyon, difüzyon, projeksiyon, kaldırma kuvveti ve alt-ızgara viskozitesi gibi --- kayıp fonksiyonuna eklenmiş cezalar olarak değil, doğrudan ağı oluşturan nöronlar olarak tanımlanmasıdır. Bir başka deyişle ``fizik kısıt değil, yapıdır'' (\textit{physics as structure, not penalty}) ilkesidir. TGV3D probleminde INNATE'in ön sonuçları PINN'e kıyasla iki büyüklük mertebesine kadar düşük hata oranları ile umut verici çıkmıştır. Ancak TGV3D, kaldırma kuvveti içermeyen, periyodik küp içinde basit yapılı bir akış geçişi problemidir; ısı transferi, alt-ızgara modellemesi ve uzun zaman ufuklarında zorlanmış-doğal konveksiyon dengesinin korunması gibi gerçekçi mühendislik gereksinimlerini içermez. INNATE'in saf-fizik mimari prensibinin gerçek anlamda doğrulanması için çok daha zorlu bir test problemine ihtiyaç duyulmuştur.

Bu çalışmada söz konusu zorlu test problemi olarak \emph{üç boyutlu karışık konveksiyon} seçilmiştir. Problem, iki farklı türbülans-üretme mekanizmasının aynı akış alanında bir arada bulunduğu bir konfigürasyon üzerine kurulmuştur:
\begin{enumerate}
  \item \textbf{Kolmogorov zorlaması}: Büyük ölçeklerde mekanik enerji girişi sağlayan, yatay yönde sinüsoidal değişen sabit-amplitüdlü bir kuvvet alanı, $k_{f}=4$ dalga sayısında \cite{kolmogorov_1941}.
  \item \textbf{Rayleigh--B\'enard tipi termal kaldırma}: Alttan ısıtılan ve üstten soğutulan kararsız sıcaklık katmanlanması altında, Boussinesq yaklaşımı ile devreye giren kaldırma kuvveti.
\end{enumerate}
Bu çalışmada ``karışık konveksiyon'' terimi, Richardson sayısı rejimine değil, yukarıdaki iki mekanizmanın birlikte aktif olduğu denklem yapısına atfen kullanılmaktadır; rejim olarak çalışma noktası ($\mathrm{Ri}\approx 1.4\times 10^{-3}$) zorlanmış konveksiyon baskın bir uç noktayı temsil eder. Akış $[0,L_{x}]\times[0,L_{y}]\times[0,L_{z}]$ periyodik küp benzeri bir bölgede tanımlanmış, $\mathrm{Re}=10\,000$ ve $\mathrm{Ra}=10^{5}$ koşullarında çözülmüştür. Bu parametre kombinasyonu, ön çalışmalarda yapılan sistematik tarama sonucunda enerji içeren büyük yapıların ortaya çıktığı ancak kaldırma kuvvetinin sayısal kararsızlık yaratmadığı bir çalışma noktası olarak belirlenmiştir.

\section{Çalışmanın Amacı}

Bu tezin temel amacı, ön dönemde TGV3D problemi üzerinden kavramsal düzeyde önerilen INNATE saf-fizik mimari yaklaşımının, bir önceki çalışmada gözlenen sınırlamaları aşacak şekilde, ısı transferi içeren, çok ölçekli ve gerçek mühendislik problemlerinin temel davranışlarını barındıran bir akış probleminde \emph{çalışma kapasitesini ortaya koymaktır}. Bu kapsamda alt amaçlar şu şekilde sıralanabilir:

\begin{enumerate}
  \item Boussinesq yaklaşımı altında üç boyutlu karışık konveksiyon denklemlerini çözmek için fiziksel operatörleri doğrudan mimari yapı olarak içeren bir saf-fizik sinir operatörü tasarlamak.
  \item Smagorinsky alt-ızgara modeli içeren çok katmanlı bir fraksiyonel adım IMEX zaman integratörünü, modelin temel öğrenilebilir omurgası olarak inşa etmek.
  \item Bu mimarinin geliştirilmesi sırasında ortaya çıkan parametre kaçışı (amplitude pumping), anti-fiziksel Nusselt değerleri ve dengesiz kayıp denkliği gibi başarısızlık modlarını sistematik biçimde belgelemek ve çözüm stratejilerini geliştirmek.
  \item Eğitilen modeli, bağımsız olarak üretilmiş düşük-Mach LES referans veri seti ile türbülans kinetik enerjisi, Nusselt sayısı, sıcaklık dalgalanma şiddeti ve enerji spektrumu metrikleri üzerinden karşılaştırarak doğrulamak.
  \item Modelin başarılı ve sınırlı olduğu yönleri akademik dürüstlük çerçevesinde ayırt etmek; özellikle termal alanın yapısal kısıtlarını açıkça ortaya koymak.
\end{enumerate}

\section{Çalışmanın Kapsamı}

Çalışmanın kapsamı bilinçli biçimde \emph{tek bir problem üzerinde derin doğrulama} biçiminde tutulmuştur. Çok sayıda problem üzerinde yüzeysel testler yapmak yerine, tek bir karışık konveksiyon konfigürasyonu seçilmiş ve bu konfigürasyon için (i) tutarlı LES referans veri seti üretimi, (ii) saf-fizik sinir operatörü tasarımı, (iii) sistematik başarısızlık analizi ve düzeltme, (iv) kapsamlı çok-metrikli doğrulama aşamaları sırayla uygulanmıştır. Bu yaklaşımla, mimarinin gerçekten ne yapıp ne yapamadığı konusunda kesin yargılara varılması hedeflenmiştir.

Çalışmanın kapsamı dışında bırakılan başlıklar şunlardır:

\begin{itemize}
  \item Periyodik olmayan, duvar sınırlı (örneğin kapalı kutu Rayleigh--B\'enard) geometriler. Bu sınırlama hem alt-ızgara modelinin tasarımını basitleştirmek hem de spektral analizleri doğrudan uygulayabilmek için tercih edilmiştir.
  \item Çoklu Reynolds--Rayleigh sayısı kombinasyonları üzerinden genelleme sınamaları. Çalışmada yalnızca $\mathrm{Re}=7000$ ve $\mathrm{Re}=10\,000$ için LES referansları üretilmiş, eğitim ve karşılaştırmalar bu iki çalışma noktası ile sınırlandırılmıştır.
  \item Doğrudan sayısal simülasyon (DNS). $\mathrm{Re}=10\,000$ için DNS gridi $\mathcal{O}(10^{9})$ noktayı aşacağından bu çalışmada uygulanabilir değildir; LES alt-ızgara modelli referans olarak benimsenmiştir.
  \item Diğer makine öğrenmesi paradigmaları (Fourier neural operator, DeepONet vb.) ile sayısal karşılaştırma. Bu yaklaşımlar literatürde \cite{li_fno_2021, lu_deeponet_2021, lesnets_2024} ele alınmıştır; karşılaştırma çerçevesi gelecek çalışmalara bırakılmıştır.
  \item Bitirme projesi 1 kapsamında geliştirilen PINN modeli ile doğrudan sayısal karşılaştırma. Bu tezde ele alınan problem, ön dönemdeki TGV3D probleminden esasen daha karmaşıktır (ısı transferi içerir, $\mathrm{Re}$ on kat daha yüksektir, kaldırma kuvveti ve alt-ızgara modeli içerir); dolayısıyla iki çalışmanın doğrudan eş zamanlı kıyaslanması anlamlı değildir. Bu tezin odağı, INNATE'in zor problemde \emph{çalışıp çalışmadığını} ortaya koymaktır.
\end{itemize}

\section{Yöntem ve Aşamalar}

Çalışma altı aşamada yürütülmüştür. (i) Problem seçimi ve parametrelendirme aşamasında, kararsızlık göstermeyen ancak yine de türbülans içeren bir karışık konveksiyon konfigürasyonu sistematik tarama ile belirlenmiştir. (ii) LES referans aşamasında, Smagorinsky alt-ızgara modeli ve Boussinesq kaldırma terimini içeren, anizotropik damping ile dengelenmiş bir LES çözücüsü ile $\mathrm{Re}=7000$ ve $\mathrm{Re}=10\,000$ için 60\,000 zaman adımı uzunluğunda referans veri setleri üretilmiştir. (iii) Saf-fizik mimari tasarım aşamasında, 20 katmanlı fraksiyonel adım IMEX integratörü, fiziksel nöron kütüphanesinin (adveksiyon, difüzyon, projeksiyon, kaldırma, alt-ızgara, zorlama) tek bir tutarlı omurga altında birleştirilmesiyle inşa edilmiş ve toplam parametre sayısı 9905 ile sınırlandırılmıştır. (iv) Ön eğitim aşamasında, model anti-fiziksel davranışlar (parametre kaçışı, $\mathrm{Nu}=200$--$430$ mertebesinde Nusselt değerleri, theta-rms LES'in 5--10 katı düzeyinde) sergilemiş; bu başarısızlıklar bölüm 4'te ayrıntılı belgelenmiştir. (v) Kademe~1+2+3 disiplinli eğitim aşamasında parametre hijyeni, PDE artık kaybı ve minimal 5-terimli kayıp seti birlikte uygulanarak modelin anti-fizik davranışları giderilmiştir. (vi) Doğrulama ve karşılaştırma aşamasında eğitilen model, bağımsız LES referansı ile türbülans kinetik enerjisi, Nusselt sayısı, sıcaklık dalgalanma şiddeti ve enerji spektrumu üzerinden karşılaştırılmıştır.

\section{Tezin Düzeni}

Tezin geri kalanı şu şekilde düzenlenmiştir. Bölüm~\ref{ch:teorik}'de problemin teorik temelleri (Boussinesq yaklaşımı, karışık konveksiyon denklemleri, türbülans istatistikleri, LES ve Smagorinsky alt-ızgara modeli, PINN'den INNATE'e evrim, saf-fizik nöron paradigması) ele alınır. Bölüm~\ref{ch:mimari}'de INNATE mimarisinin ayrıntıları (fiziksel nöron kütüphanesi, fraksiyonel adım IMEX integratörü, Spectral-Cs v3 parametrizasyonu, Kademe~1+2+3 eğitim disiplini, kayıp fonksiyonu tasarımı ve eğitim hattı) sunulur. Bölüm~\ref{ch:bulgular}'da LES referans üretimi, eğitim sürecinde ortaya çıkan başarısızlık modlarının sistematik analizi, Kademe~1+2+3 ile elde edilen iyileştirmeler ve modelin LES referansına karşı çok-metrikli karşılaştırması verilir. Bölüm~\ref{ch:sonuc}'de ise INNATE yaklaşımının karışık konveksiyon problemindeki başarı ve sınırları tartışılır, akademik dürüstlük çerçevesinde modelin yapısal kısıtları kayıt altına alınır ve gelecek çalışmalara yönelik öneriler sunulur.
